\subsection{双重直接极限。直接极限的乘积}

设 I, L 为两个有向集，并设 \(I \times L\) 为它们的乘积（§ 1，第 4 号），该乘积在乘积预序下仍是一个有向集。考虑一个集合直接系统 \((\mathrm{E}_{\alpha}^{\lambda}, f_{\beta\alpha}^{\mu\lambda})\) （相对于 \(I \times L\) ）. 于是我们得到

\[
f _ {\gamma \alpha} ^ {\nu \lambda} = f _ {\gamma \beta} ^ {\mu \nu} \circ f _ {\beta \alpha} ^ {\mu \lambda}\tag{29}
\]

每当 \(\alpha \leqslant \beta \leqslant \gamma\) 并且 \(\lambda \leqslant \mu \leqslant \nu\) .

让 \(\mathbf{E}\) 或 \(\lim_{\overrightarrow{\alpha, \lambda}} \mathbf{E}_{\alpha}^{\lambda}\) 表示此直接系统的直接极限。对每个 \(\lambda \in L\) ，令 \(g_{\beta \alpha}^{\lambda} = f_{\beta \alpha}^{\lambda \lambda}: \mathbf{E}_{\alpha}^{\lambda} \to \mathbf{E}_{\beta}^{\lambda}\) 。于是由(29)我们有

\[
g _ {\gamma \alpha} ^ {\lambda} = g _ {\gamma \beta} ^ {\lambda} \circ g _ {\beta \alpha} ^ {\lambda}\tag{30}
\]

\[
\text { 当 } \alpha \leqslant \beta \leqslant \gamma ;
\]

换言之， \((\mathbf{E}_{\alpha}^{\lambda}, g_{\beta \alpha}^{\lambda})\) 是相对于I的一个集合直接系统。设 \(\mathbf{F}^{\lambda} = \lim_{\overrightarrow{\alpha}} \mathbf{E}_{\alpha}^{\lambda}\) 表示其直接极限。若 \(\lambda \leqslant \mu\) 是L的固定元素，则由(29)可知映射 \(h_{\alpha}^{\mu \lambda} = f_{\alpha \alpha}^{\lambda \mu}: \mathbf{E}_{\alpha}^{\lambda} \to \mathbf{E}_{\alpha}^{\mu}\) 构成一个映射直接系统。设 \(h^{\mu \lambda} = \varinjlim_{\alpha} h_{\alpha}^{\mu \lambda}: \mathbf{F}^{\lambda} \to \mathbf{F}^{\mu}\) 表示该映射系统的直接极限。

若 \(\lambda \leqslant \mu \leqslant \nu\) 在 \(\mathbf{L}\) 中，那么

\[
h ^ {\nu \lambda} = h ^ {\nu \mu} \circ h ^ {\mu \lambda}\tag{31}
\]

（第6号，命题6，推论2），因此 \((\mathbf{F}^{\lambda}, h^{\mu \lambda})\) 是相对于 \(\mathbf{L}\) 的一个集合直接系统。设 \(\mathbf{F} = \lim \mathbf{F}^{\lambda}\) 为其直接极限。我们将定义一个典范双射 \(\mathbf{E} \to \mathbf{F}\) 。为此，令 \(g_{\alpha}^{\lambda}\) 表示典范映射 \(\mathbf{E}_{\alpha}^{\lambda} \to \mathbf{F}^{\lambda}\) ， \(h^{\lambda}\) 表示典范映射 \(\mathbf{F}^{\lambda} \to \mathbf{F}\) ，并令 \(u_{\alpha}^{\lambda} = h^{\lambda} \circ g_{\alpha}^{\lambda}\) 。如果 \(\alpha \leqslant \beta\) 并且 \(\lambda \leqslant \mu\) ，那么我们有

\[
\begin{array}{r l} u _ {\beta} ^ {\mu} \circ f _ {\beta \alpha} ^ {u \lambda} & = h ^ {\mu} \circ g _ {\beta} ^ {\mu} \circ f _ {\beta \alpha} ^ {u \lambda} = h ^ {\mu} \circ g _ {\beta} ^ {\mu} \circ f _ {\beta \alpha} ^ {\mu \mu} \circ f _ {\alpha \alpha} ^ {\mu \lambda} = h ^ {\mu} \circ g _ {\alpha} ^ {\mu} \circ f _ {\alpha \alpha} ^ {\mu \lambda} \\ & = h ^ {\mu} \circ h ^ {\mu \lambda} \circ g _ {\alpha} ^ {\lambda} = h ^ {\lambda} \circ g _ {\alpha} ^ {\lambda} = u _ {\lambda} ^ {\lambda} \end{array}
\]

由(29)及 \(h^{\mu \lambda}\) 的定义。因此， \(u_{\alpha}^{\lambda}\) 构成一个相对于 \(\mathbf{I} \times \mathbf{L}\) 的映射直接系统。令 \(u = \lim_{\overrightarrow{\alpha, \lambda}} u_{\alpha}^{\lambda}: \mathbf{E} \to \mathbf{F}\) 。我们

将证明 \(u\) 是双射的，通过应用第6号命题6的判据。首先， \(\mathbf{F}\) 是诸集合 \(h^{\lambda}(\mathbf{F}^{\lambda})\) 的并集，以及每一个 \(\mathbf{F}^{\lambda}\) 是诸集合 \(g_{\alpha}^{\lambda}(\mathbf{E}_{\alpha}^{\lambda})\) 的并集；因此F是诸集合的并

\[
h ^ {\lambda} (g _ {\alpha} ^ {\lambda} (\mathbf {E} _ {\alpha} ^ {\lambda})) = u _ {\alpha} ^ {\lambda} (\mathbf {E} _ {\alpha} ^ {\lambda}).
\]

接下来，令 \(x, y\) 为 \(\mathbf{E}_{\alpha}^{\lambda}\) 的两个元素，使得 \(u_{\alpha}^{\lambda}(x) = u_{\alpha}^{\lambda}(y)\) ，即 \(h^{\lambda}(g_{\alpha}^{\lambda}(x)) = h^{\lambda}(g_{\alpha}^{\lambda}(y))\) 。那么（第5号，引理1）存在 \(\mu \geqslant \lambda\) 使得 \(h^{\mu \lambda}(g_{\alpha}^{\lambda}(x)) = h^{\mu \lambda}(g_{\alpha}^{\lambda}(y))\) ，即 \(g_{\alpha}^{\mu}(f_{\alpha \alpha}^{\mu \lambda}(x)) = g_{\alpha}^{\mu}(f_{\alpha \alpha}^{\mu \lambda}(y))\) ；同样地，存在 \(\beta \geqslant \alpha\) 使得 \(g_{\beta \alpha}^{\mu}(f_{\alpha \alpha}^{\mu \lambda}(x)) = g_{\beta \alpha}^{\mu}(f_{\alpha \alpha}^{\mu \lambda}(y))\) （第5号，引理1），即 \(f_{\beta \alpha}^{\mu \lambda}(x) = f_{\beta \alpha}^{\mu \lambda}(y)\) ；并且这表明（第6号，命题6） \(u\) 是单射的。因此我们证明了：

命题9。若 \((\mathbf{E}_{\alpha}^{\lambda}, f_{\beta \alpha}^{\mu \lambda})\) 是相对于两个有向集的乘积 \(\mathbf{I} \times \mathbf{L}\) 的一个集合直接系统，则（在典范双射意义下）我们有

\[
\lim _ {\overrightarrow {\alpha}, \overrightarrow {\lambda}} \mathbf {E} _ {\alpha} ^ {\lambda} = \lim _ {\overrightarrow {\lambda}} (\lim _ {\overrightarrow {\alpha}} \mathbf {E} _ {\alpha} ^ {\lambda}).\tag{32}
\]

推论。设 \((\mathbf{E}_{\alpha}^{\prime \lambda}, f_{\beta \alpha}^{\prime u_{\lambda}})\) 是另一个相对于 \(\mathbf{I} \times \mathbf{L}\) 的集合直接系统，并且对于每个 \((\alpha, \lambda) \in \mathbf{I} \times \mathbf{L}\) 让 \(u_{\alpha}^{\lambda}\) 为从 \(\mathbf{E}_{\alpha}^{\lambda}\) 到 \(\mathbf{E}_{\alpha}^{\prime \lambda}\) 中的一个映射，使得诸 \(u_{\alpha}^{\lambda}\) 构成一个映射直接系统。则我们得到

\[
\varinjlim_ {\alpha , \lambda} u _ {\alpha} ^ {\lambda} = \varinjlim_ {\lambda} (\varinjlim_ {\alpha} u _ {\alpha} ^ {\lambda}).\tag{33}
\]

我们将其验证留给读者。

命题10。设 \((\mathbf{E}_{\alpha}, f_{\beta \alpha})\) 并且 \((\mathbf{E}_{\alpha}', f_{\beta \alpha}')\) 是两个集合直接系统，均相对于同一个有向集合 I。设 \(\mathbf{E} = \varinjlim \mathbf{E}_{\alpha}\) , \(\mathbf{E}' = \varinjlim \mathbf{E}_{\alpha}'\) ，并且让 \(f_{\alpha}: \mathbf{E}_{\alpha} \to \mathbf{E}\) , \(f_{\alpha}': \mathbf{E}_{\alpha}' \to \mathbf{E}'\) 表示典范映射，对每个 \(\alpha \in \mathbf{I}\) 。那么 \((\mathbf{E}_{\alpha} \times \mathbf{E}_{\alpha}', f_{\beta \alpha} \times f_{\beta \alpha}')\) 是一个集合直接系统， \((f_{\alpha} \times f_{\alpha}')\) 是一个映射直接系统，并且 \(\varinjlim (f_{\alpha} \times f_{\alpha}')\) 是一个双射

\[
\varinjlim (\mathbf {E} _ {\alpha} \times \mathbf {E} _ {\alpha} ^ {\prime}) \rightarrow (\varinjlim \mathbf {E} _ {\alpha}) \times (\varinjlim \mathbf {E} _ {\alpha} ^ {\prime}).\tag{34}
\]

命题的前两个断言立即得到验证。为证明 \(g = \lim (f_{\alpha} \times f_{\alpha}')\) 是双射的，我们应用第6号命题6。显然 \(\mathbf{E} \times \overrightarrow{\mathbf{E}'}\) 是诸集合 \(f_{\alpha}(\mathbf{E}_{\alpha}) \times f_{\alpha}'(\mathbf{E}_{\alpha}')\) 的并集；因此 \(g\) 是满射。如果 \((x, x')\) 并且 \((y, y')\) 是 \(\mathbf{E}_{\alpha} \times \mathbf{E}_{\alpha}'\) 的两个元素，使得 \(f_{\alpha}(x) = f_{\alpha}(y)\) 并且 \(f_{\alpha}'(x') = f_{\alpha}'(y')\) ，那么（第5号，引理1）存在两个元素 \(\beta, \gamma\) （\(\mathbf{I}\) 的元素）使得 \(\beta \geqslant \alpha, \gamma \geqslant \alpha\) ，以及 \(f_{\beta \alpha}(x) = f_{\beta \alpha}(y), f_{\gamma \alpha}'(x') = f_{\gamma \alpha}'(y')\) ；由于 \(\mathbf{I}\) 是有向的，存在 \(\delta \in \mathbf{I}\) 使得 \(\delta \geqslant \beta\) 并且 \(\delta \geqslant \gamma\) ；因此 \(f_{\delta \alpha}(x) = f_{\delta \alpha}(y)\) 并且 \(f_{\delta \alpha}'(x') = f_{\delta \alpha}'(y')\) 。这就完成了证明。

双射g称为典范的。

推论。设 \((\mathrm{F}_{\alpha}, g_{\beta\alpha})\) 和 \((\mathrm{F}_{\alpha}^{\prime}, g_{\beta\alpha}^{\prime})\) 是两个相对于I的集合直接系统，且对每个 \(\alpha \in I\) 让 \(u_{\alpha}: E_{\alpha} \to F_{\alpha}, u_{\alpha}^{\prime}: E_{\alpha}^{\prime} \to F_{\alpha}^{\prime}\) 为映射，使得 \((u_{\alpha})\) 和 \((u_{\alpha}^{\prime})\) 是两个映射直接系统。则 \((u_{\alpha} \times u_{\alpha}^{\prime})\) 是一个映射的直接系统，并且（在典范双射的意义下）我们有

\[
\lim _ {\rightarrow} \left(u _ {\alpha} \times u _ {\alpha} ^ {\prime}\right) = (\lim _ {\rightarrow} u _ {\alpha}) \times (\lim _ {\rightarrow} u _ {\alpha} ^ {\prime}).\tag{35}
\]

我们将其验证留给读者。

